\begin{lemma}[Part IV, Algorithm (B) and invariant (16)]
Fix the following Case 1 data.  Let
\(1_1<\cdots<1_{p+1}\) be a strictly increasing finite list of exact
PE slots, all supported at \(v_n\); in the extension-base application
this is the exact class \(L_1\) and \(1_{p+1}=n\).  The chosen missing color at \(u\) is
\(\gamma_{1_{b+1}}\), no earlier gamma in this list is missing at
\(u\), and
\(L_1^*=\{1_1,\ldots,1_b\}\).  Assume the exact gamma, delta, stored-pair,
PE-kind, and supporting-vertex fields, the chain identities
\(\gamma_{1_{a+2}}=\delta_{1_{a+1}}\), and their stated distinctness.
Assume that every pair indexed by \(L_1^*\) is strongly interchangeable
under \(\varrho_1\), that \(T_j\) is an ordered initial segment of each
corresponding \(T_t\), and that (15) holds at \(v_n\).  If
\(\varrho_1\) is proper and leaves exactly \(e_0\) uncolored, then there
is an \noderef{ExtensionBaseAlgorithmBTrace} satisfying
\noderef{ExtensionBaseAlgorithmBInvariant}.

Equivalently, Algorithm (B) can be carried out in decreasing order on
the exact set \(L_1^*\), and before every iteration \(B(t)\):
\(\delta_t\) is missing at \(u\); the switched path contains at most the
single vertex \(v_n\) of \(T_t\), with \(v_n\) not an end; and every
unprocessed pair remains strongly interchangeable after the switch.
\end{lemma}
\begin{proof}
Put
\[
 A=\noderef{ExtensionBaseCaseOneIndices}(\mathtt{indices},b),
 \qquad
 U=\bigcup_{t\in A}S_t .                                     \tag{1}
\]
Thus \(A=L_1^*\), and the terminal hypothesis is
\(\overline{\varrho_1}(v_n)\cap U=\varnothing\).  If \(b=0\), the
definition of \(A\) makes it empty.  Take the zero-step
\noderef{ExtensionBaseAlgorithmBTrace}, with sole coloring
\(\varrho_1\).  Its properness and exact-one-uncolored clauses are the
hypotheses, and every per-step or absent-slot clause mentioning a member
of \(A\) is vacuous.  This gives
\noderef{ExtensionBaseAlgorithmBInvariant}.  Henceforth let
\(m=b.\mathtt{val}>0\).

For \(0\leq a<m\), set
\[
 r_a=m-1-a,\qquad t_a=\mathtt{indices}(r_a).                  \tag{2}
\]
Here and below a natural number in the displayed range is regarded as its
unique element of \(\operatorname{Fin}(p+1)\).  The list
\[
 t_0=\mathtt{indices}(m-1),\ t_1=\mathtt{indices}(m-2),
 \ldots,\ t_{m-1}=\mathtt{indices}(0)                         \tag{3}
\]
has image exactly \(A\).  If \(a<c<m\), then \(r_c<r_a<b\), so strict
increase of \(\mathtt{indices}\) gives \(t_c<t_a\).  Define
\(\varrho^0=\varrho_1\), and recursively define
\[
 K_a=\noderef{TwoColorComponent}
   (G,\varrho^a,\gamma_{r_a},\delta_{r_a},\mathsf{inl}(u)),
 \qquad
 \varrho^{a+1}=\noderef{TwoColorComponentSwap}
   (G,\varrho^a,\gamma_{r_a},\delta_{r_a},\mathsf{inl}(u)).   \tag{4}
\]
At step \(a\), use position \(r_a\), slot \(t_a\), colors
\(\gamma_{r_a},\delta_{r_a}\), support \(v_n\), and coloring
\(\varrho^a\).  Equations (2)--(4), the exact support, kind, gamma, and
delta hypotheses, and the image and order facts above supply every field
of an \noderef{ExtensionBaseAlgorithmBTrace}; call it \(\mathcal B\).

We first record precisely what the color-chain hypotheses say.  Define
\[
 c_0=\gamma_0,\qquad c_{q+1}=\delta_q\quad(0\leq q\leq p).    \tag{5}
\]
The chain identity gives \(\gamma_q=c_q\) for every \(q\): this is the
definition for \(q=0\), and for \(q>0\) it is
\(\gamma_q=\delta_{q-1}=c_q\).  The map in (5) is exactly
\(\operatorname{Fin.cons}(\gamma_0,\delta)\), so its assumed
injectivity says that the colors \(c_0,\ldots,c_{p+1}\) are pairwise
distinct.  In particular,
\[
 \gamma_q=c_q\ne c_{q+1}=\delta_q,                            \tag{6}
\]
and, whenever \(q<r\), \(\gamma_q=c_q\) differs from both colors
\(c_r,c_{r+1}\) switched at position \(r\).  The exact stored-pair
hypothesis also gives
\[
        \gamma_q,\delta_q\in U\qquad(q<b).                   \tag{7}
\]

We prove the invariant for \(\mathcal B\) by induction down (3).  The
induction assertion immediately before position \(r=r_a\) consists of
the following facts:
\[
\begin{array}{ll}
({\rm I}) & \varrho^a\text{ is proper and leaves exactly }e_0
            \text{ uncolored};\\
({\rm II})& \delta_r\in\overline{\varrho^a}(u),\text{ and every }
            \gamma_q\ (q\leq r)\text{ is present at }u;\\
({\rm III})&\text{the pairs at the current and all lower unprocessed
            positions are strongly interchangeable under }\varrho^a;\\
({\rm IV})&\text{every color in }U\text{ is present at }v_n
            \text{ under }\varrho^a.
\end{array}                                                     \tag{8}
\]
For \(a=0\), (I) is the pair of coloring hypotheses.  Since
\(r_0=m-1\), the chain identity gives
\(\delta_{r_0}=\gamma_b\); this is missing at \(u\) by the selected-color
hypothesis.  Each \(\gamma_q\) with \(q\leq r_0<b\) is nonmissing by
minimality.  By \noderef{MissingColors}, nonmissing means membership in
\noderef{PresentColors}, so (II) follows.  Assertion (III) is the initial
strong-interchangeability hypothesis.  Finally, if \(c\in U\), then the
terminal hypothesis says \(c\notin\overline{\varrho_1}(v_n)\); unfolding
\noderef{MissingColors} and \noderef{PresentColors} produces an edge
identity \(e\) with \noderef{Incident}\((G,v_n,e)\) and
\(\varrho_1(e)=\mathsf{some}(c)\).  This is (IV).

Fix a step \(a\), write \(r=r_a\),
\(\alpha=\gamma_r\), \(\beta=\delta_r\),
\(\varphi=\varrho^a\), and let \(K=K_a\).  We now verify every per-step
clause and then re-establish (8) for the next position.

By (II), \(\beta\) is missing and \(\alpha\) is present at \(u\).
The definitions \noderef{MissingColors} and \noderef{PresentColors}
therefore give no edge incident with \(u\) having value
\(\mathsf{some}(\beta)\), and give an incident edge \(e_\alpha\) with
\(\varphi(e_\alpha)=\mathsf{some}(\alpha)\).  This edge is unique:
two distinct such identities share \(u\), hence are adjacent by
\noderef{AdjacentEdges}, and then contradict properness as defined in
\noderef{IsProperEdgeColoring}.  By \noderef{TwoColorIncidenceGraph},
the edge-node \(\mathsf{inr}(e_\alpha)\) is adjacent to
\(\mathsf{inl}(u)\), so it lies in the rooted component \(K\).  Moreover
it is the only neighbor of \(u\) in that incidence graph: a neighbor is
an incident \(\alpha\)- or \(\beta\)-colored edge, and the preceding two
sentences respectively give uniqueness and nonexistence.  Thus \(K\)
contains a genuine edge-node and \(u\) has a singleton, in particular
subsingleton, neighbor set.  Distinctness (6), together with the
spanning path supplied by \noderef{TwoColorComponentSpanningPath}, now
gives exactly the \noderef{TwoColorNontrivialPathComponent} certificate required
in invariant (16); the spanning path has full support \(K\), and \(u\)
is one of its ends.

The current pair is strongly interchangeable by (III).  Unfold
\noderef{ETTStrongInterchangeable}, and suppose that
\(x\in V(T_{t_a})\), \(x\ne v_n\), and
\(\mathsf{inl}(x)\in K\).  Because \(x\) and \(u\) lie in the same
connected component, the rooted-component definition
\noderef{TwoColorComponent} gives
\[
 \operatorname{Comp}_{\alpha,\beta}(x)=K.                    \tag{9}
\]
Strong interchangeability gives a
\noderef{TwoColorComponentCycle} whose support is the entire left side
of (9), not merely a cycle contained in it.  Hence its
\noderef{TwoColorCycle} contains every node of \(K\), including \(u\).
The two directions of a simple closed cycle at \(u\) give two distinct
neighbors of \(u\).  This contradicts the subsingleton neighbor set
proved above from the rooted path endpoint.  Therefore
\[
 \mathsf{inl}(x)\in K\ \text{ and }\ x\in V(T_{t_a})
       \quad\Longrightarrow\quad x=v_n.                      \tag{10}
\]
This is the exact tree-intersection clause of the invariant.

For completeness, the corresponding assertion for paths used at earlier
iterations really does descend to the present tree; it is not an
unsupported comparison of two \(T_t\)'s.  If \(c<a\), (3) gives
\(t_a<t_c\).  The ETT retains its certified series by
\noderef{ExtendedTashkinovTree}, and \noderef{TashkinovSeries} supplies a
\noderef{TashkinovSeriesStep} certificate at every positive intermediate
slot.  Each noninitial SE, PE, or RE certificate has an ordered-extension
field.  Composing these fields from \(t_a\) through \(t_c\), by composing
the length inequalities and entrywise equalities in
\noderef{TreeSequenceInitialSegment}, proves that \(T_{t_a}\) is an
ordered initial segment of \(T_{t_c}\).  The definition
\noderef{TreeSequenceVertexSet} then gives
\(V(T_{t_a})\subseteq V(T_{t_c})\).  The already proved intersection
bound for the earlier path \(K_c\) on \(T_{t_c}\) consequently restricts
to the same bound on \(T_{t_a}\), exactly as used in the manuscript's
decreasing-order argument.

We next prove both the terminal endpoint clause and preservation of the
strong terminal assertion (IV).  By (7) and (IV), both \(\alpha\) and
\(\beta\) are present at \(v_n\).  Thus \noderef{PresentColors} gives
incident identities \(e_\alpha^n,e_\beta^n\) of the two respective
colors.  They are distinct by (6), since one coloring function cannot
give two different \(\mathsf{some}\)-values to the same edge.  If
\(\mathsf{inl}(v_n)\in K\), then
\noderef{TwoColorIncidenceGraph} makes both edge-nodes adjacent to
\(v_n\), hence members of \(K\).  They are two distinct neighbors.
The bound \noderef{TwoColorIncidenceDegreeAtMostTwo}, applied using
properness in (I), shows that they exhaust the neighbor set at \(v_n\).
In particular that neighbor set is not subsingleton, which proves that
\(v_n\) is not an end, as required by invariant (16).

Now let \(c\in U\), and choose from (IV) an incident \(c\)-colored edge at
\(v_n\).  If \(v_n\notin K\), no incident edge at \(v_n\) can be changed:
any changed edge has color \(\alpha\) or \(\beta\), so its incidence with
\(v_n\) would put the vertex-node in \(K\) by
\noderef{TwoColorIncidenceGraph} and
\noderef{TwoColorComponent}.  If \(v_n\in K\), the preceding paragraph
shows that the swap exchanges the colors on exactly the two incident
component edges \(e_\alpha^n,e_\beta^n\).  Hence both current colors
remain present, on the opposite identities.  If
\(c\notin\{\alpha,\beta\}\), inspection of
\noderef{TwoColorComponentSwap} shows that the chosen \(c\)-colored edge
is fixed whether or not its edge-node belongs to \(K\).  These cases prove
that every \(c\in U\) remains present at \(v_n\) under
\(\varrho^{a+1}\), establishing (IV) for the next stage.  Notice that
the use of an other color here is legitimate: (5)--(6), namely the
injectivity of \(\operatorname{Fin.cons}(\gamma_0,\delta)\), identifies
and separates all color values occurring in the stored union.

We now establish the missing-color descent at \(u\) without suppressing
the incidence argument.  The only incident \(\alpha\)-edge under
\(\varphi\) is \(e_\alpha\), it belongs to \(K\), and the swap changes its
value to \(\mathsf{some}(\beta)\).  There was no incident \(\beta\)-edge.
Consequently an incident edge could have color \(\alpha\) after the swap
only by either being an old \(\alpha\)-edge outside \(K\), excluded by
uniqueness and membership of \(e_\alpha\), or being an old \(\beta\)-edge
inside \(K\), excluded by \(\beta\)'s missingness.  Thus \(\alpha\) is
missing and \(\beta\) is present at \(u\) after the swap, by
\noderef{MissingColors}, \noderef{PresentColors}, and
\noderef{Incident}.  For every \(q<r\), (5)--(6) says that
\(\gamma_q=c_q\) differs from both switched colors \(c_r,c_{r+1}\).
The definition \noderef{TwoColorComponentSwap} therefore fixes the truth
of \(\varphi(e)=\mathsf{some}(\gamma_q)\) for every edge \(e\), so the
incident witness making \(\gamma_q\) present at \(u\) survives.  If a
next position exists, its index is \(r-1\), and the chain identity is
\[
              \delta_{r-1}=\gamma_r=\alpha .                 \tag{11}
\]
The newly missing \(\alpha\), (11), and the surviving lower gammas prove
(II) for that next position.

Properness after the switch follows from (6), (I), and
\noderef{TwoColorComponentSwapProper}.  We also verify the unique
uncolored edge pointwise.  For every edge identity \(e\), expansion of
\noderef{TwoColorComponentSwap} gives
\[
 \varrho^{a+1}(e)=\mathsf{none}
       \quad\Longleftrightarrow\quad
 \varrho^a(e)=\mathsf{none}.                                 \tag{12}
\]
Outside \(K\) the value is fixed; inside \(K\), the two
\(\mathsf{some}\)-values are exchanged and every other value, including
\(\mathsf{none}\), is fixed.  Unfolding
\noderef{IsOneEdgeUncolored}, its first clause and (12) give
\(\varrho^{a+1}(e_0)=\mathsf{none}\).  For \(e\ne e_0\), its second
clause supplies \(\varrho^a(e)=\mathsf{some}(d)\); the swap either fixes
that value or exchanges it with the other current
\(\mathsf{some}\)-value, so \(\varrho^{a+1}(e)\) is again colored.
This discharges both clauses and completes (I) for the next stage.

It remains to propagate strong interchangeability.  Take an arbitrary
unprocessed \(s\in A\) after \(t_a\), and arbitrary colors
\(\gamma_s,\delta_s\) satisfying the invariant's exact gamma- and
delta-field premises at \(s\).  The exact enumeration (2)--(3) writes
\(s=\mathtt{indices}(q)\) for a unique \(q<r\); hence strict increase gives
\(s<t_a\).  Comparing the supplied field equalities with the hypotheses
at position \(q\), and using injectivity of \(\mathsf{some}\), gives
\(\gamma_s=\gamma_q\) and \(\delta_s=\delta_q\).  The current and
survivor slots both have exact PE kind and support \(v_n\), by the
trace fields at \(t_a\) and the corresponding hypotheses at \(q\).
Thus every exact support, kind, gamma, delta, and index premise has now
been instantiated.  Their two pairs are strongly interchangeable by
(III), the current rooted component is the path component proved above,
and (I) supplies the exact properness and one-uncolored hypotheses.
Therefore the same-support, earlier-slot alternative of
\noderef{StrongInterchangeabilityKempePreserved}, instantiated with these
exact supports, fields, and the inequality \(s<t_a\), proves that the
pair at \(s\) remains strongly interchangeable under
\(\varrho^{a+1}\).  This holds for every unprocessed \(s\), so it proves
(III) for the next stage.

Finally, the supplied hypothesis \(T_j\subseteq T_{t_a}\) is an ordered
initial-segment statement at every processed position.  Together with
(I), the missing-delta assertion in (II), the path certificate, (10),
the terminal nonendpoint argument, and the survivor argument, it gives
every per-step clause of
\noderef{ExtensionBaseAlgorithmBInvariant}.  The induction also gives
properness and exact-one-uncoloredness for the final coloring.  Since the
trace enumerates all of \(A\), a slot in \(A\) cannot be absent from every
trace position; hence the invariant's final absent-slot survivor clause is
vacuous.  All clauses of the invariant now hold for \(\mathcal B\), which
is the required trace.
\end{proof}
